\subsection{معماری و پیاده‌سازی}
کلاس \lr{FLAPCompressor} در فایل \lr{Codes/framework/compression\_engine/flap.py} پیاده‌سازی شده و تحت شناسه‌های \lr{flap} و \lr{flap\_compressor} در \lr{COMPRESSOR\_REGISTRY} ثبت گردیده است. ویژگی‌های معماری و مهندسی این ماژول عبارتند از:
\begin{itemize}
    \item \textbf{سازگاری کامل با \lr{Conv1D} و \lr{Linear}:} پیاده‌سازی با تشخیص خودکار نوع لایه، ترانهاده مناسب وزن‌ها را در لایه‌های \lr{c\_attn} و \lr{c\_proj} در \lr{GPT-2} اعمال کرده و سپس ماسک تنکی را اعمال می‌نماید.
    \item \textbf{پشتیبانی از تنکی سخت‌افزاری $2:4$ (\lr{Structured 2:4 Pruning}):} امکان فعال‌سازی پرچم \lr{structured\_24} برای استفاده از شتاب‌دهنده‌های تانسوری \lr{NVIDIA Ampere} و معماری‌های جدیدتر جهت افزایش توان عملیاتی استنتاج.
    \item \textbf{حفاظت از لایه‌های حساس (\lr{Critical Layer Protection}):} فیلترهای استخراج لایه‌ها از دستکاری طبقه‌بند نهایی (\lr{classifier})، سر مدل زبان (\lr{lm\_head})، نرمال‌سازی لایه‌ها و پروجکشن‌های لورای موجود ممانعت می‌کنند.
    \item \textbf{سازوکار بازگشت ایمن (\lr{Safe Fallback}):} در صورت تخطی خطای بازسازی فروبنیوس از آستانه مجاز (\lr{max\_error=0.35})، لایه بدون تغییر باقی مانده و مشخصات آن در گزارش فشرده‌سازی ثبت می‌شود.
\end{itemize}
